\chapter{Dashboard}
\label{ch:dashboard}

\section{Arquitectura}

El dashboard del proyecto es un \emph{sitio web estático} generado por
\texttt{dashboard/html/build.py} y servido por un pequeño backend Flask
(\texttt{dashboard/html/server.py}) que añade unos cuantos endpoints REST
sobre el DWH y los artefactos serializados. El build genera un único
fichero \texttt{index.html} (de aproximadamente $1{,}1$\,MB) con todas las
vistas embebidas y todo el JavaScript de interactividad inline; los datos
en vivo se obtienen vía \emph{fetch} contra los endpoints del servidor. El
servidor también detecta automáticamente si PostgreSQL está accesible y
conmuta entre dos modos:

\begin{itemize}
  \item \textbf{Modo \emph{vivo}}: las vistas de \emph{Visión General} y
        \emph{Productos} consultan directamente \texttt{dwh.fact\_sales}
        y \texttt{dwh.fact\_returns} para mostrar KPIs del periodo,
        evolución diaria y \emph{rankings} de tiendas y productos.
  \item \textbf{Modo \emph{offline}}: ante la ausencia de PostgreSQL,
        las vistas dependientes de DWH se reemplazan por su versión
        construida sólo desde los \emph{parquet} ya generados.
\end{itemize}

\section{Endpoints REST}

El backend expone los siguientes endpoints, todos en JSON:

\begin{lstlisting}[language=Python,
  caption={Endpoints del servidor del dashboard, en \texttt{server.py}.},
  label={lst:endpoints}]
GET /api/health                       # ping; n_customers, dwh online?
GET /api/customers/search             # busqueda + filtros sobre customer_segments
GET /api/customer/<customer_id>       # ficha completa: comparativas + radar + timeline
GET /api/cluster/<cluster_id>         # perfil estadistico + top clientes + top productos
GET /api/cluster/predict              # simulador: feature_vec -> cluster + (pc1,pc2)
GET /api/cluster/feature_ranges       # rangos (p01..p99) por feature para los sliders
\end{lstlisting}

\noindent
La búsqueda admite filtros por \texttt{cluster}, mínimo y máximo de CLTV,
mínimo y máximo de frecuencia, y orden por CLTV, recencia, frecuencia o
nombre. El simulador \texttt{/api/cluster/predict} carga los tres
artefactos \texttt{scaler/pca/kmeans.joblib} en memoria y aplica
exactamente la misma cadena utilizada en el entrenamiento, de manera que
las predicciones sobre clientes ficticios sean reproducibles.

\section{Las seis vistas}

La navegación del dashboard se ofrece como un sumario con seis capítulos,
en línea con la estética editorial del informe. Cada capítulo del
dashboard corresponde a un módulo bajo
\texttt{dashboard/html/views/}\,:

\begin{enumerate}
  \item \textbf{Visión General} (\texttt{overview.py}): KPIs del año,
        inventario verificado del DWH, evolución diaria con media móvil
        a $7$ días, mezcla por categoría, \emph{treemap} categoría/marca,
        estacionalidad y top 10 de tiendas
        (figura~\ref{fig:dash-overview}).
  \item \textbf{Clientes} (\texttt{clientes.py}): hallazgo principal
        (porcentaje de \emph{one-shot}), distribución del CLTV con
        histograma logarítmico, ECDF y densidad por cluster, curva de
        Lorenz con coeficiente de Gini, y top $20$ de clientes
        (figura~\ref{fig:dash-clientes}).
  \item \textbf{Segmentación} (\texttt{segmentacion.py}): tres tarjetas
        de cluster con KPI por cabeza, mapa PCA con \emph{hexbin},
        tamaño de cluster, varianza explicada, \emph{heatmap}
        $z$-score y simulador interactivo
        (figura~\ref{fig:dash-segmentacion}).
  \item \textbf{Productos} (\texttt{productos.py}): KPIs del periodo,
        top $15$ de productos por ingresos, descomposición de
        devoluciones por motivo, top productos por valor devuelto
        (figura~\ref{fig:dash-productos}).
  \item \textbf{Ficha Cliente} (\texttt{ficha.py}): buscador en vivo
        sobre $5\,750$ clientes con filtros por cluster, CLTV o
        frecuencia; al seleccionar, despliega cabecera, comparativas,
        radar de percentiles e historial de compras leído del DWH
        (figura~\ref{fig:dash-ficha}).
  \item \textbf{Metodología} (\texttt{metodologia.py}): mapa de
        KPI~$\to$~fórmula~$\to$~origen, listado del pipeline ETL, la
        propia justificación del $k=3$, varianza explicada y estadísticos
        clave (figura~\ref{fig:dash-metodologia}).
\end{enumerate}

\section{Capturas de pantalla}

Las capturas de las seis vistas se obtuvieron con el servidor Flask en
ejecución sobre \texttt{localhost:8001} y exportadas en alta resolución.

\begin{figure}[H]
  \centering
  \includegraphics[width=\linewidth]{dashboard_overview.png}
  \caption[Dashboard, Visión General]{Dashboard, vista
  \emph{Visión General}: KPIs del año en curso, inventario verificado del
  DWH y comienzo de la serie diaria de ingresos.}
  \label{fig:dash-overview}
\end{figure}

\begin{figure}[H]
  \centering
  \includegraphics[width=\linewidth]{dashboard_clientes.png}
  \caption[Dashboard, Clientes]{Dashboard, vista \emph{Clientes y CLTV}:
  hallazgo principal (un único \emph{ticket} en el $87\%$ de la base),
  histograma logarítmico del CLTV, curva de Lorenz con
  $\text{Gini}=0{,}907$ y comparación mediana--media por cluster.}
  \label{fig:dash-clientes}
\end{figure}

\begin{figure}[H]
  \centering
  \includegraphics[width=\linewidth]{dashboard_segmentacion.png}
  \caption[Dashboard, Segmentación]{Dashboard, vista
  \emph{Segmentación}: tarjetas-resumen por cluster, mapa PCA con
  contornos, justificación del $k=3$, varianza explicada y simulador
  interactivo del modelo.}
  \label{fig:dash-segmentacion}
\end{figure}

\begin{figure}[H]
  \centering
  \includegraphics[width=\linewidth]{dashboard_productos.png}
  \caption[Dashboard, Productos]{Dashboard, vista
  \emph{Productos y devoluciones}: KPIs del periodo, top $15$ por
  ingresos, distribución por motivo de devolución y mapa de hechos
  utilizados (\texttt{fact\_sales}, \texttt{fact\_returns},
  \texttt{dim\_product}, \texttt{dim\_return\_reason}).}
  \label{fig:dash-productos}
\end{figure}

\begin{figure}[H]
  \centering
  \includegraphics[width=\linewidth]{dashboard_ficha.png}
  \caption[Dashboard, Ficha Cliente]{Dashboard, vista
  \emph{Ficha Cliente}: buscador en vivo y panel detallado del cliente
  con radar de percentiles, tabla comparativa contra mediana del cluster
  y mediana global, e historial de compras agregado por categoría.}
  \label{fig:dash-ficha}
\end{figure}

\begin{figure}[H]
  \centering
  \includegraphics[width=\linewidth]{dashboard_metodologia.png}
  \caption[Dashboard, Metodología]{Dashboard, vista
  \emph{Metodología}: pipeline ETL, justificación del $k=3$ por curva del
  codo, varianza explicada y estadísticos clave del CLTV.}
  \label{fig:dash-metodologia}
\end{figure}

\section{Decisiones de diseño}

El dashboard se ha construido siguiendo una estética \emph{editorial} o
de \emph{data-storytelling} (próxima al estilo de los \emph{longreads} de
publicaciones como Pudding o Bloomberg) antes que un patrón clásico de
``rejilla de KPIs''. Las decisiones más relevantes:

\paragraph{Una columna principal con párrafos cortos.}
Cada gráfico se introduce con una pequeña descripción que orienta la
lectura, lo que evita que el lector tenga que reconstruir el contexto
visual por su cuenta.

\paragraph{Paletas de colores reservadas a la rentabilidad.}
El color de los gráficos no codifica categorías arbitrarias sino
métricas semánticamente comparables: el margen aparece en rampa
azul\,$\to$\,verde, la concentración (Lorenz) en azul, y los clusters en
una terna estable (azul, ámbar, rojo) que se repite en todas las vistas
de segmentación.

\paragraph{Modo \emph{vivo} vs. \emph{offline} con misma maqueta.}
El sitio detecta la disponibilidad del DWH al arrancar y conmuta los
componentes pertinentes sin alterar la maqueta, lo que permite usarlo
como prototipo de presentación incluso cuando no hay base de datos
disponible.

\paragraph{Plotly inline + Flask para el detalle.}
Los gráficos se generan en el servidor mediante \texttt{plotly.io.to\_html}
y se embeben directamente en \texttt{index.html}, lo que evita
dependencias en tiempo de carga; los componentes verdaderamente
interactivos (\emph{ficha de cliente}, \emph{simulador}, \emph{cluster
detail}) son los que consumen los endpoints REST.
